\section{Introduction}\label{sec:intro}

\subsection{Critical collapse}

In 1993 Choptuik studied numerically the gravitational collapse of a spherically symmetric massless scalar field
and found, at the threshold between dispersion and black hole formation, a universal solution with three striking
features~\cite{Choptuik1993}: it is approached, in all the one-parameter families of initial data that were evolved, when the
parameter is tuned to the threshold; it is discretely self-similar, reproducing itself on smaller and smaller scales with an echoing period
$\Delta\approx3.44$; and near the threshold the mass of the black hole scales like $(p-p_*)^\gamma$ with a universal
exponent $\gamma\approx0.37$. These \emph{critical phenomena} have since been found in many matter models, and are
reviewed in~\cite{GundlachMartinGarcia2007}. For the spherically symmetric Einstein--scalar field system, Christodoulou proved that
sufficiently small data disperse and sufficiently large data form trapped surfaces~\cite{Christodoulou1986,
Christodoulou1991,Christodoulou1993}.

The standard explanation is dynamical~\cite{KoikeHaraAdachi1995,Gundlach1997}. The critical solution is a fixed
point (in the discretely self-similar case, a periodic orbit) of the evolution in logarithmic time, with a stable
manifold of codimension one, which is the threshold itself, and a \emph{single} unstable mode, growing like
$e^{\lambda T}$. A near-critical solution stays close to the critical one until the unstable mode has grown to a
fixed size, and dimensional analysis then gives $\gamma=1/\lambda$, with a periodic modulation in the discretely
self-similar case. Gundlach computed the perturbations of Choptuik's solution numerically and found exactly one growing mode,
with $\gamma=0.374\pm0.001$~\cite{Gundlach1997}, and Mart\'in-Garc\'ia and Gundlach found numerically that all
nonspherical perturbations decay~\cite{MartinGarciaGundlach1999}. A necessary ingredient of this picture is that the
critical solution has exactly one unstable mode; it is this spectral statement that we prove, for spherically
symmetric perturbations.

\subsection{Existence, and the open spectral question}

The existence of Choptuik's spacetime was proved by Reiterer and Trubowitz~\cite{RT2019}, with a computer-assisted
argument in the tradition of Lanford's proof of the Feigenbaum conjectures~\cite{Lanford1982}. They construct a
real analytic, discretely self-similar solution on a neighbourhood of the past light cone of the singularity, as
the fixed point of a contraction near high-precision numerical data, and determine its echoing constant and a
second invariant to $80$ digits. About perturbations they write: ``We do not study perturbations about the
Choptuik spacetime, that are relevant for many of the (conjectured) properties of the solution. Our paper can be
the starting point for a rigorous investigation of this kind. This would be interesting, because current
numerical results seem to be inconclusive''~\cite[\S1]{RT2019}.

This paper carries out such an investigation for spherically symmetric perturbations.

\subsection{The result}

We work with the solution of~\cite{RT2019}, in their coordinates $(\tau,\xi)$ and variables $\omega$, in which the
self-similarity is the translation $\tau\mapsto\tau+2\pi$. Linear perturbations of Floquet type,
$e^{s\tau}$ times a periodic profile, are governed by a family of operators $L(s)$, the linearization of the
Reiterer--Trubowitz system, together with a linearized constraint map $C(s)$. The exponent $s$ corresponds to the
growth rate $\lambda=2\pi s/K$ in the usual logarithmic time, where $e^{-K}$ is the factor by which the
self-similarity rescales lengths.

\begin{introthm}[informal version of the Main Theorem, \cref{sec:main}]
Under spherically symmetric linear perturbations that are smooth in time and analytic in the radial variable up to
and across the light cone, and modulo gauge, Choptuik's spacetime has exactly one unstable mode. The mode is real and algebraically simple, and its growth rate satisfies
\[
  \lambda\in[2.674040,\ 2.674115],\qquad 1/\lambda\in[0.373955,\ 0.373966].
\]
More precisely, per Floquet class the operator $L$ has exactly four roots in the half plane $\Re s>0$, each simple
and with a real representative: a gauge mode, which is the translation of the time of the singularity; two roots whose eigenvectors
violate the constraints; and the physical unstable mode. On the imaginary axis, $L$ has only the root $s=0$, the
translation of the phase. If the light cone of the singularity is held fixed, the gauge mode is no longer removed by the
smaller group of gauge transformations, and the count becomes two; the additional class corresponds to translating
the time of the singularity, not to a second instability.
\end{introthm}

The value of $1/\lambda$ agrees with the measured and computed critical exponent; this is an external check, since
the relation $\gamma=1/\lambda$ concerns the nonlinear evolution. The theorem is spectral: it locates the roots of
$L$ and identifies the physical ones, but says nothing about the linearized evolution semigroup or the nonlinear
dynamics. It concerns spherically symmetric perturbations only, in the regularity class stated above. These
limitations are discussed in \cref{sec:discussion}.

\subsection{Strategy}

The proof has three parts.

\emph{From the geometric problem to the operator} (\cref{sec:physical}). A physical mode is a solution of the
linearized Einstein--scalar field equations of Floquet type, and only its class modulo gauge transformations is
meaningful. We show that every such mode with $\Re s>0$ can be put in the form of~\cite{RT2019} by a gauge
transformation, and then corresponds to an element of $\ker L(s)\cap\ker C(s)$; that conversely every such element
comes from a physical mode; and that the only gauge modes with $\Re s>0$ are the multiples of one explicit mode at
$s=\mu$, the infinitesimal translation of the singular point. The gauge completeness is global on the domain
of~\cite{RT2019} and uses that the light cone lies inside it. The physical multiplicity of a root is then a matter
of linear algebra on its root space.

\emph{Counting} (\cref{sec:counting}). The roots of $L$ with large real part are excluded by a bound on the
numerical range, computed in a function space in which the background acts pointwise. In the remaining bounded
region, the roots are counted with an operator version of Rouch\'e's theorem due to Gohberg and
Sigal~\cite{GohbergSigal1971}: after a rank-two deflation of the two gauge roots, $L$ is compared along the
boundary with a block-diagonal reference whose zeros are the eigenvalues of a $14016\times14016$ matrix, and these
are counted from a Schur decomposition with a certified residual.

\emph{Identification} (\cref{sec:roots}). Each root is enclosed, and shown to be algebraically simple, by a bordered
Newton--Kantorovich argument. The eigenvectors at two of them are shown to violate the constraints by explicit lower
bounds. For the remaining root, the constraints cannot be checked on an approximate eigenvector; instead, an
identity $KC=ML$ propagates the constraints, and a separate certificate shows that the constraint propagation
operator $K$ is injective there.

All the computations are done in exact rational arithmetic, in ball arithmetic, or as a posteriori bounds on
floating-point computations. They take a few days on ordinary computers.

\subsection{What is new}

The tools are classical; what is new is their combination for this operator, and the results. The paper turns into a
theorem a spectral picture that was known only numerically: it gives a rigorous count of the unstable modes of
Choptuik's spacetime; a rigorous enclosure of the
physical growth rate; the classification of all the roots with positive real part; and the passage from the geometric
problem to the operator, including the completeness of the gauge. The certificates, a verification program that
re-checks them, the data and a reproduction guide are publicly available (\cref{sec:computation-data}).

\subsection{Organization}

\Cref{sec:setting} recalls the formulation of~\cite{RT2019} and defines $L$, $C$ and $K$. \Cref{sec:main} states the
Main Theorem. \Cref{sec:physical} relates physical modes to the operator, \cref{sec:counting} counts the roots and
\cref{sec:roots} identifies them and completes the proof. \Cref{sec:computation} describes the computations, their
verification and the use of AI assistants, and \cref{sec:discussion} discusses the limitations and the nonspherical
problem. The appendices contain the proofs of the structural lemmas, explicit formulas and the list of certificates.
